\section{R ile çözümlü uygulama}
\label{sec:r-b01}

\textbf{Soru.} Bölümdeki beş öğrencilik öğretim örneğini R'de oluşturun.
Gözlem birimini, değişken türlerini, program frekanslarını ve ortalama başarıyı
belirleyin. Bu küçük veri gerçek bir araştırmadan alınmış değildir.

\begin{lstlisting}[style=prismR]
veri <- data.frame(
  devam_saati = c(8, 10, 7, 12, 9),
  basari = c(68, 75, 64, 83, 72),
  program = factor(c("A", "A", "B", "B", "A"))
)
str(veri)
summary(veri)
table(veri$program)
prop.table(table(veri$program))
colSums(is.na(veri))
c(ogrenci = nrow(veri),
  ortalama_devam = mean(veri$devam_saati),
  ortalama_basari = mean(veri$basari))
\end{lstlisting}

\begin{outputguidebox}
\textbf{Çözüm ve kontrol değerleri.} Her satır bir öğrencidir; veri
5 satır ve 3 değişken içerir. \texttt{devam\_saati} ve \texttt{basari}
nicel, \texttt{program} nominal kategoriktir. A programında 3, B programında
2 öğrenci vardır. Ortalama devam 9,2 saat, ortalama başarı 72,4 puandır;
hiçbir sütunda eksik değer yoktur.
\end{outputguidebox}

\textbf{Yorum.} \texttt{factor} program kodlarının kategori olarak
saklanmasını sağlar. Bir program etiketinin ortalamasını almak anlamlı
değildir. Beş satırın özetlenmesi, devamın başarıya neden olduğunu veya
öğrencilerin bir evreni temsil ettiğini göstermez.

\textbf{Pekiştirme ve yanıt.} Programların örneklemdeki paylarını
\texttt{prop.table} çıktısından okuyun: A için 0,60, B için 0,40 elde edilir. Bunlar evren oranı değil örneklem oranlarıdır.